\documentclass[12pt,a4paper]{article}
\usepackage{fontspec}
\usepackage{graphicx}
\usepackage{booktabs}
\usepackage[a4paper,left=26mm,right=26mm,top=29mm,bottom=27mm,headheight=25pt]{geometry}
\usepackage{xcolor}
\usepackage{fancyhdr}
\usepackage[hidelinks,unicode]{hyperref}
\usepackage{xepersian}

\settextfont[
  Path=../assets/,
  Scale=1.12,
  BoldFont=B-NAZANIN.TTF
]{B-NAZANIN.TTF}

\graphicspath{{assets/}}
\definecolor{BookletAccent}{HTML}{573626}
\definecolor{BookletMuted}{HTML}{59514D}
\definecolor{BookletRule}{HTML}{C0AEA3}

\linespread{1.22}
\setlength{\parindent}{0pt}
\setlength{\parskip}{0.55em plus 0.15em}
\setlength{\abovecaptionskip}{9pt}
\setlength{\belowcaptionskip}{6pt}
\renewcommand{\arraystretch}{1.22}
\raggedbottom

\makeatletter
\renewcommand{\section}{%
  \@startsection{section}{1}{0pt}%
    {-3.0ex plus -0.8ex minus -0.2ex}{1.1ex plus 0.2ex}%
    {\normalfont\Large\color{black}}}
\renewcommand{\subsection}{%
  \@startsection{subsection}{2}{0pt}%
    {-2.5ex plus -0.6ex minus -0.2ex}{0.8ex plus 0.2ex}%
    {\normalfont\large\color{black}}}
\renewcommand*{\l@section}[2]{\@dottedtocline{1}{0em}{1.8em}{#1}{#2}}
\makeatother

\pagestyle{fancy}
\fancyhf{}
\fancyhead[R]{\small جزوه‌های فنی چرم‌دوزی}
\fancyfoot[L]{\small\thepage}
\renewcommand{\headrulewidth}{0pt}
\renewcommand{\footrulewidth}{0pt}

\newcommand{\BookletCover}[2]{%
  \hypersetup{pdftitle={#1}}%
  \begin{titlepage}
    \vspace*{22mm}%
    {\color{BookletAccent}\large جزوه‌های فنی چرم‌دوزی\par}
    \vspace{5mm}%
    {\color{BookletRule}\rule{\linewidth}{0.8pt}\par}
    \vspace{25mm}%
    {\color{black}\fontsize{30}{37}\selectfont #1\par}
    \vspace{5mm}%
    {\color{BookletMuted}\Large #2\par}
  \end{titlepage}%
}


\begin{document}
\BookletCover{نمونۀ قالب جزوه}{چیدمان متن، جدول، شکل و منابع}
\tableofcontents
\clearpage

\section{ثبت اندازه‌های یک نمونه}
عنوان هر ستون مشخص می‌کند که عددهای آن ستون چه کمیتی را نشان می‌دهند و با چه واحدی ثبت شده‌اند. در جدول~\ref{tab:measurements}، هر اندازه به یکی از قسمت‌های مشخص‌شده در شکل~\ref{fig:parts} تعلق دارد.

\subsection{معرفی اصطلاح فنی}
اصطلاح \lr{Full-grain leather} به چرمی اشاره دارد که لایۀ طبیعی رویه را به‌طور کامل حفظ کرده است \cite{lwg}.

\begin{table}[htbp]
  \centering
  \caption{اندازه‌های فرضی دو قسمت از یک قطعه}
  \label{tab:measurements}
  \begin{tabular}{rcc}
    \toprule
    قسمت & طول (میلی‌متر) & عرض (میلی‌متر) \\
    \midrule
    الف & ۴۵ & ۲۰ \\
    ب   & ۳۰ & ۲۰ \\
    \bottomrule
  \end{tabular}
\end{table}

\begin{figure}[htbp]
  \centering
  \setlength{\fboxsep}{0pt}
  \setlength{\fboxrule}{0.6pt}
  \fbox{\parbox[c][20mm][c]{45mm}{\centering الف}}%
  \hspace{3mm}%
  \fbox{\parbox[c][20mm][c]{30mm}{\centering ب}}
  \caption{طرح شماتیک قسمت‌های اندازه‌گیری‌شده}
  \label{fig:parts}
\end{figure}

\begin{thebibliography}{9}
  \bibitem{lwg}
  \lr{Leather Working Group. \emph{\href{https://www.leatherworkinggroup.com/fileadmin/uploads/lwg/Knowledge/LWG_List_of_Definitions_v1.1.pdf}{List of Definitions}}. Version 1.1, January 2024.}\newline
  مدخل «\mbox{\lr{Full grain}}».
\end{thebibliography}
\end{document}
